\chapter{Anexos}
\label{chap:anexos}

\section{Extracto de la Hoja de Datos del Transistor MOSFET 2N7000}

El dispositivo activo utilizado en la totalidad de las experiencias del trabajo práctico es un transistor MOSFET de
canal N tipo enriquecimiento en encapsulado plástico estándar TO-92, fabricado por ON Semiconductor. En el
Cuadro~\ref{tab:datasheet_ratings_tp5} y en el Cuadro~\ref{tab:datasheet_caract_tp5} se sintetizan los valores límites
absolutos de operación y las características eléctricas de pequeña señal más representativas obtenidas de su hoja de datos.

\begin{table}[!ht]
  \centering
  \resizebox{\textwidth}{!}{
  \begin{tabular}{|l|c|c|c|}
    \hline
    \textbf{Parámetro Máximo Admisible} & \textbf{Símbolo} & \textbf{Valor Límite} & \textbf{Unidad} \\
    \hline
    Tensión drenador-surtidor & $V_{DSS}$ & $60$ & $\V$ \\
    Tensión drenador-compuerta ($R_{GS} \le 1.0\Mohm$) & $V_{DGR}$ & $60$ & $\V$ \\
    Tensión compuerta-surtidor continua & $V_{GSS}$ & $\pm 20$ & $\V$ \\
    Tensión compuerta-surtidor pulsada ($t_p \le 50\us$) & $V_{GSM}$ & $\pm 40$ & $\V$ \\
    Corriente continua de drenador ($T_A = 25^\circ\text{C}$) & $I_D$ & $200$ & $\mA$ \\
    Corriente pulsada de drenador & $I_{DM}$ & $500$ & $\mA$ \\
    Disipación total de potencia a $T_A = 25^\circ\text{C}$ & $P_D$ & $350$ & $\mW$ \\
    Rango de temperatura de juntura operativa & $T_J$ & $-55 \text{ a } +150$ & $^\circ\text{C}$ \\
    \hline
  \end{tabular}
  }
  \caption{Límites máximos absolutos de operación del MOSFET 2N7000 (ON Semiconductor).}
  \label{tab:datasheet_ratings_tp5}
\end{table}

\begin{table}[!ht]
  \centering
  \resizebox{\textwidth}{!}{
  \begin{tabular}{|l|c|c|c|c|c|}
    \hline
    \textbf{Característica Eléctrica} & \textbf{Símbolo} & \textbf{Mín.} & \textbf{Típ.} & \textbf{Máx.} & \textbf{Unidad} \\
    \hline
    Tensión umbral de compuerta ($V_{DS} = V_{GS}, I_D = 1\mA$) & $V_{GS(\mathrm{th})}$ & $0.8$ & $2.1$ & $3.0$ & $\V$ \\
    Corriente de fuga de drenador con compuerta a masa ($V_{DS} = 48\V, V_{GS} = 0$) & $I_{DSS}$ & -- & -- & $1.0$ & $\uA$ \\
    Corriente de fuga compuerta-cuerpo ($V_{GS} = \pm 15\V, V_{DS} = 0$) & $I_{GSS}$ & -- & -- & $\pm 100$ & $\mathrm{nA}$ \\
    Transconductancia directa ($V_{DS} = 10\V, I_D = 200\mA$) & $g_{fs}$ & $100$ & $320$ & -- & $\mS$ \\
    Resistencia estática drenador-surtidor en encendido ($V_{GS} = 10\V, I_D = 500\mA$) & $R_{DS(\mathrm{on})}$ & -- & $1.2$ & $5.0$ & $\ohm$ \\
    Capacitancia de entrada ($V_{DS} = 25\V, V_{GS} = 0, f = 1\MHz$) & $C_{iss}$ & -- & $20$ & $60$ & $\pF$ \\
    Capacitancia de salida ($V_{DS} = 25\V, V_{GS} = 0, f = 1\MHz$) & $C_{oss}$ & -- & $11$ & $25$ & $\pF$ \\
    Capacitancia de transferencia reversa ($V_{DS} = 25\V, V_{GS} = 0, f = 1\MHz$) & $C_{rss}$ & -- & $4.0$ & $5.0$ & $\pF$ \\
    \hline
  \end{tabular}
  }
  \caption{Características eléctricas estáticas y dinámicas del MOSFET 2N7000.}
  \label{tab:datasheet_caract_tp5}
\end{table}

La tensión umbral obtenida experimentalmente en el laboratorio ($V_{th} = 1.60104\V$) se encuentra perfectamente dentro
del rango admisible de fabricación ($0.8\V \text{ a } 3.0\V$). Asimismo, la potencia disipada por el transistor en reposo es:
\begin{equation}
  P_{D,Q} = V_{DSQ} \cdot I_{DQ} = 6.390\V \cdot 1.0206\mA \approx 6.52\mW
\end{equation}
magnitud que representa apenas el $1.86\%$ de la potencia máxima disipable ($350\mW$), asegurando que el dispositivo
opera frío y con una confiabilidad térmica sobresaliente.

\section{Netlist Completo de Simulación en LTspice}

A continuación se transcribe el netlist circuital íntegro correspondiente al archivo esquemático
\texttt{ej\_Amp.SurtidorComun1.asc}:

\begin{lstlisting}[style=ltspice, caption={Netlist de simulación para el amplificador en surtidor común en LTspice.}, label={list:netlist_surtidor_comun}]
* Amplificador en Configuracion Surtidor Comun con MOSFET 2N7000
* Archivo fuente: sim/ej_Amp.SurtidorComun1.asc
M1 d g s s 2N7000
VDD n001 0 15
RD n001 d 6.8k
RS s 0 1.5k
R1 g 0 130k
R2 n001 g 470k
C1 in g 100uF
C2 s 0 100uF
C3 d out 100uF
RL out 0 10k
Vi in 0 SINE(0 45mV 1kHz)

* Subcircuito de fabricante para el transistor 2N7000
.SUBCKT 2N7000 1 2 3
M1 9 7 8 8 MM L=100u W=100u
.MODEL MM NMOS LEVEL=1 IS=1e-32 VTO=1.60104 LAMBDA=0 KP=0.053905
+ CGSO=1.79115e-07 CGDO=1.0724e-11
RS 8 3 1.10523
D1 3 1 MD
.MODEL MD D IS=2.71011e-10 RS=0.0140826 N=1.5 BV=60 IBV=1e-05
+ EG=1.16084 XTI=3.00131 TT=0 CJO=3.41211e-11 VJ=4.67429 M=0.899864
RDS 3 1 2.4e+11
RD 9 1 0.0001
RG 2 7 2.18034
D2 4 5 MD1
.MODEL MD1 D IS=1e-32 N=50 CJO=7.93181e-11 VJ=0.643298 M=0.9
D3 0 5 MD2
.MODEL MD2 D IS=1e-10 N=0.400165 RS=3.00002e-06
RL 5 10 1
FI2 7 9 VFI2 -1
VFI2 4 0 0
EV16 10 0 9 7 1
CAP 11 10 1.58786e-10
FI1 7 9 VFI1 -1
VFI1 11 6 0
RCAP 6 10 1
D4 0 6 MD3
.MODEL MD3 D IS=1e-10 N=0.400165 RS=0 BV=1000 IBV=1mA
.ENDS 2N7000

.tran 50m
.op
.backanno
.end
\end{lstlisting}
